\section{Plane clipping and its verification}
\label{app:clip}

This appendix sets out the region construction of
Section~\ref{sec:methods:clip} in enough detail to be reimplemented, derives
the closed-form quantities used to verify it, and states the tolerance
treatment on which its robustness depends.

\subsection{The algorithm}
\label{app:clip:alg}

Let $(\mathcal{V},\mathcal{F})$ be a closed, outward-oriented triangle mesh and
let the cutting plane be $x = c$. We construct the closed mesh bounding the
solid $\{x > c\}$; the complement and the slab of the neck model follow by
reversing the sense and by clipping twice.

Each vertex is classified by its signed distance $d_i = x_i - c$ into one of
three states,
\begin{equation}
\mathrm{code}(i) =
\begin{cases}
+1, & d_i > \tau, \\
\phantom{+}0, & |d_i| \le \tau, \\
-1, & d_i < -\tau,
\end{cases}
\label{eq:appclassify}
\end{equation}
with $\tau$ the tolerance of Section~\ref{app:clip:tol}. Vertices in state $0$
are snapped exactly onto the plane, $x_i \leftarrow c$, and are thereafter
treated as belonging to the kept side.

A facet is discarded if all three codes are $\le 0$ and retained unchanged if
all three are $\ge 0$. Otherwise it is clipped. Traversing its three edges in
order, the vertex $i$ is emitted whenever $\mathrm{code}(i) \ge 0$, and an
intersection point is emitted whenever an edge changes sign \emph{strictly},
that is when $\mathrm{code}(i)\,\mathrm{code}(j) < 0$, at
\begin{equation}
\mathbf{p} = \mathbf{v}_i + t\,(\mathbf{v}_j - \mathbf{v}_i),
\qquad
t = \frac{d_i}{d_i - d_j},
\label{eq:appintersect}
\end{equation}
with the $x$-coordinate of $\mathbf{p}$ set to $c$ exactly rather than left to
rounding. The requirement that the sign change be strict is what prevents a
vertex lying on the plane from generating two coincident intersection points;
this is developed in Section~\ref{app:clip:tol}. The resulting polygon is
emitted as a fan from its first vertex after consecutive coincident points have
been removed, and is discarded if fewer than three distinct vertices remain.

Intersection points are cached by the edge that produced them. An edge is
identified by the unordered pair of its endpoint indices, so the two facets
sharing a cut edge receive the same new vertex index. Created per facet
instead, they would receive two distinct indices at the same position: the
volume would be unchanged, since it is a sum of signed tetrahedra, but the
surface would be split along the entire cut and would fail the manifold test of
Section~\ref{app:clip:verify}.

\subsection{Closing the opening}
\label{app:clip:cap}

Every edge of the retained surface whose two endpoints both lie on the plane
and which has non-zero length is a boundary edge. These are traced into closed
loops by following, from each edge, the unique unvisited outgoing edge at its
head; the traversal terminates when it returns to its start or runs out. Each
loop is then capped separately, to its own apex: for a loop $\mathcal{L}$ with
vertex set $\mathcal{V}_\mathcal{L}$, the apex $\mathbf{a}$ is the centroid of
$\mathcal{V}_\mathcal{L}$ projected onto the plane, and for each
$(i,j) \in \mathcal{L}$ the triangle $(\mathbf{a}, \mathbf{v}_j, \mathbf{v}_i)$
is added, the reversed traversal orienting the cap outward.

That a fan closes a planar loop correctly, for any coplanar apex, follows from
the signed-area identity: for a closed loop $\{\mathbf{v}_k\}$ in a plane with
unit normal $\nhat$,
\begin{equation}
\sum_k \tfrac{1}{2}\,\nhat \cdot
\big[(\mathbf{v}_k - \mathbf{a}) \times (\mathbf{v}_{k+1} - \mathbf{a})\big]
\;=\; A_{\mathrm{signed}},
\label{eq:appfan}
\end{equation}
independently of $\mathbf{a}$, provided $\mathbf{a}$ lies in the plane.

The identity holds for each loop separately, and this is why the loops must be
traced rather than fanned together. A single apex shared between two disjoint
loops still gives the correct total signed area---so the volume would be
right---but the resulting triangles are not a surface: the two loops become
connected through a vertex that belongs to both, and the mesh is no longer a
manifold. The condition arises whenever a cross-section is multiply connected
or the region falls into more than one piece, and both occur in this study. The
Arrokoth shape model is supplied as two unjoined shells, so a plane beyond its
neck produces a region in two pieces, each with its own cap.

\subsection{Tolerance and degenerate facets}
\label{app:clip:tol}

The tolerance is set to
\begin{equation}
\tau = 10^{-10} \, \lVert \mathbf{v}_{\max} - \mathbf{v}_{\min} \rVert ,
\label{eq:apptol}
\end{equation}
the bounding-box diagonal of the mesh, so that it scales with the body.

Its purpose is to prevent a specific degeneracy. Suppose a vertex $k$ lies
exactly on the plane, $d_k = 0$, and is classified by a strict test as
\emph{outside}. Both edges incident on it from inside then satisfy
$d_i > 0 > d_k$ formally, and Eq.~\eqref{eq:appintersect} returns $t = 1$ for
each, placing both intersection points at $\mathbf{v}_k$ itself. The clipped
polygon acquires two coincident vertices, and the triangles formed with them
have zero area, on the surface and on the cap alike.

The consequence is more severe than a small geometric error. A facet of zero
area has an undefined normal, and in the Werner--Scheeres formulation the
per-facet contribution is assembled from that normal; the evaluation returns
NaN not at points near the facet but at every field point on the body. The
failure is therefore total rather than local, and it is parameter-dependent: at
$\xsplit = 0.1608$~km on Itokawa no such vertex is encountered and the
computation is clean, while at $0.1300$~km four are, and all $49{,}152$ field
values are lost.

Two guards are applied after clipping. Facets with a repeated vertex index are
removed, as are facets whose area falls below $10^{-12}$ times the median
positive area. In practice the classification of Eq.~\eqref{eq:appclassify}
prevents such facets from being created at all: on the runs reported in this
paper the number removed is zero.

Slivers---facets of small but non-zero area, produced where a vertex lies just
outside the tolerance---are not removed, and the polyhedral gravity
implementation flags them as potentially unstable. Their contribution is
proportional to their area and should vanish, but catastrophic cancellation is
conceivable when a facet millimetres across is evaluated at coordinates of
order kilometres. This was therefore measured directly. Raising $\tau$ by four
orders of
magnitude, to $10^{-6}$ of the diagonal, removes them: on 67P the minimum facet
area rose from $2.381\times10^{-12}$ to $3.660\times10^{-8}$~km$^2$ and eight
facets were eliminated. Every reported quantity agreed to the digits printed,
the dispersion at seven planes being identical to eight decimals and the
recovered contrast, improvement and excess mass unchanged in the seventh
significant figure or beyond. On a synthetic body the agreement was exact.

\subsection{Closed-form quantities for verification}
\label{app:clip:analytic}

Two solids with elementary closed forms are used to verify the construction.

For a sphere of radius $R$ cut at $x = c$, write $h = R - c$ for the height of
the resulting cap. Integrating the disc of radius $\sqrt{R^2 - x^2}$ over
$x \in [c, R]$ gives
\begin{equation}
V_{\mathrm{cap}} = \pi \int_{c}^{R} (R^2 - x^2)\,\mathrm{d}x
= \frac{\pi h^{2} (3R - h)}{3},
\qquad
\bar{x}_{\mathrm{cap}} = \frac{3\,(2R - h)^{2}}{4\,(3R - h)} .
\label{eq:appcap}
\end{equation}
For $R = 1$ and $c = 0.5$ these are $0.6544985$ and $0.6750000$.

For the slab $|x| \le a$ of the same sphere,
\begin{equation}
V_{\mathrm{slab}} = \pi \int_{-a}^{a} (R^2 - x^2)\,\mathrm{d}x
= 2\pi\!\left( R^{2} a - \frac{a^{3}}{3} \right),
\label{eq:appslab}
\end{equation}
which for $R = 1$ and $a = 0.2$ gives $1.2398819$.

The corresponding radial cone---the solid produced by the construction of
Appendix~\ref{app:cone}, whose apex is the sphere's centre and whose base is
the spherical cap, and which for a sphere is a spherical sector---has
\begin{equation}
V_{\mathrm{sector}} = \tfrac{1}{3}\,R\,A_{\mathrm{cap}}
= \tfrac{1}{3}\,R\,(2\pi R h)
= \frac{2\pi R^{2} h}{3},
\qquad
\bar{x}_{\mathrm{sector}} = \tfrac{3}{4}\!\left( R - \frac{h}{2} \right),
\label{eq:appcone}
\end{equation}
giving $1.0471976$ and $0.5625000$ for the same cut. The centroid follows from
the fan of tetrahedra sharing the apex, whose centroids lie three quarters of
the way from apex to base, the base being a spherical cap of mean
$x$-coordinate $R - h/2$.

\subsection{Verification}
\label{app:clip:verify}

Table~\ref{tab:appverify} reports what the implementation returns on an
icosphere of $20{,}480$ facets.

\begin{table}[htbp]
\centering
\caption{Verification of the clipping construction against
Eqs.~\eqref{eq:appcap}--\eqref{eq:appcone}, on a unit icosphere of $20{,}480$
facets. The residuals near $0.1\%$ are the polyhedral approximation of the
sphere: at the same subdivision the whole-body volume differs from $4\pi/3$ by
$0.05\%$. The clipping itself is exact for the given polyhedron.}
\label{tab:appverify}
\begin{tabular}{lrrr}
\toprule
Quantity & Measured & Closed form & Difference \\
\midrule
Cap volume, $c = 0.5$        & $0.653931$ & $0.654498$ & $-0.087\%$ \\
Cap centroid, $c = 0.5$      & $0.674936$ & $0.675000$ & $-0.010\%$ \\
Slab volume, $a = 0.2$       & $1.239437$ & $1.239882$ & $-0.036\%$ \\
Sector volume, $c = 0.5$     & $1.047453$ & $1.047198$ & $+0.024\%$ \\
Sector centroid, $c = 0.5$   & $0.562213$ & $0.562500$ & $-0.051\%$ \\
\midrule
Complement closure error     & \multicolumn{3}{r}{$2.0\times10^{-14}\%$} \\
Predicate violation, lobe    & \multicolumn{3}{r}{$0$} \\
Predicate violation, sector  & \multicolumn{3}{r}{$0.5$} \\
Euler characteristic, lobe   & \multicolumn{3}{r}{$2$} \\
Facets failing containment   & \multicolumn{3}{r}{$0$ (one component)} \\
\bottomrule
\end{tabular}
\end{table}

Four checks are applied to every region, and a fifth where the shape model
has more than one component. Each inspects a property the others do not, and
the set was arrived at by failure: each was added after something passed all
the checks then in place.

The \emph{analytic} comparison tests magnitude but would be passed by any
construction returning the right volume for the wrong reason.

The \emph{complement} test, $V(x>c) + V(x<c) = \Vtot$, tests internal
consistency and is passed by the sector as readily as by the lobe, since both
partition the tetrahedral sum.

The \emph{predicate} test inspects the identity of the solid, requiring
$\min_i x_i \ge c$ over the vertices of the region. It is the only one of the
first three that the sector fails.

The \emph{manifold} test inspects topology. Every directed edge must appear
exactly once with its reverse present, and the Euler characteristic
$\chi = V - E + F$ must equal $2$ for each connected component. This is
insensitive to none of the same failures: a surface split along the cut, as
described in Section~\ref{app:clip:alg}, passes the analytic, complement and
predicate tests, because none of them consults connectivity, and fails only
this one. Conversely a sector passes the manifold test and fails the predicate.

The \emph{containment} test inspects the embedding. Where the mesh has more
than one component, each facet centroid is tested by ray casting for
containment in every other component, and those found inside are excluded from
the area-weighted sums. It catches what none of the other four can: a mesh in
two closed pieces is topologically sound whatever their arrangement in space,
so a facet buried inside the opposite shell passes every previous check while
contributing a geopotential evaluated at a point that is not on the surface.
On the Arrokoth model $455$ facets, $0.943\%$ of the area, fail it, and
excluding them moves the assumed bulk density at which the recovered contrast
crosses unity from about $345$ to just below
$\SI{250}{\kilo\gram\per\cubic\metre}$
(Section~\ref{sec:nulls:arrokoth}).

The criterion is stated per connected component rather than as $\chi = 2$,
because a region cut from a body supplied as several unjoined shells is
legitimately in several pieces. Of the shape models used here, Arrokoth has
$\chi = 4$ over two components as supplied, and the 67P SPC~2017 model is not
closed at all, with one repeated and two unmatched directed edges. Both defects
are inherited by any region cut from those models, and neither affects a volume
or a potential.

The same containment machinery bears on the volume. Two components that
overlapped in space would have the material in the overlap counted twice,
since the divergence-theorem volume sums over all facets without regard to
whether they bound the same region. Measured by Monte Carlo over the
overlapping bounding box with $200{,}000$ samples, the intersection of the two
Arrokoth shells is $0.044 \pm 0.001$ per cent of the total volume, consistent
with zero: the shells touch but do not interpenetrate, and no material is
double counted. It is the surface average, not the volume, that the buried
facets corrupt.

A stress case is worth recording. On an icosphere subdivided so that $128$
vertices lie exactly on $x = 0$, cutting at that plane returns a hemisphere of
volume $2.093260$ against the closed-form $2\pi/3 = 2.094395$, a difference of
$-0.054\%$, with $\chi = 2$, zero degenerate facets and a minimum facet area of
$1.92\times10^{-4}$. A strict-inequality classification would produce
coincident points on every one of those $128$ vertices, which is why the
tolerance-based classification of Section~\ref{app:clip:tol} is required.

\subsection{Detecting the dividing plane}
\label{app:clip:neck}

The saddle detector of Section~\ref{sec:methods:neck} operates on the binned
profile $w(x)$, the ninetieth percentile of facet-centroid distance from the
$x$ axis within each of sixty bins. Two aspects of its implementation are worth
stating.

Candidates are identified on \emph{runs} of equal value rather than at single
points. A run $[j,k]$ with $w_j = \dots = w_k$ within tolerance is a local
minimum if $w_{j-1} > w_j$ and $w_{k+1} > w_k$, and its position is taken as
the midpoint. A strict point test, requiring $w_j < w_{j\pm1}$, fails on a
symmetric profile, where two samples straddle the minimum with equal values;
applied to Itokawa it returned no neck at all. It also fails on any plateau.

The prominence of Eq.~\eqref{eq:prominence} is measured against the
\emph{lower} of the two flanking maxima, in the manner of the topographic
quantity of the same name. This has a consequence worth noting: Eros scores
higher than Itokawa, $0.216$ against $0.170$, despite its shallower saddle,
because its two flanking humps are nearly equal ($7.9072$ and $7.8902$~km,
differing by $0.2\%$) while Itokawa's differ by $15\%$ ($0.1437$ and
$0.1215$~km). Prominence measures how well separated the lobes are, not how
deep the neck is, and the two are not the same.

The detector is also mesh-dependent. On Itokawa it returns $0.1702$~km on the
$20{,}000$-facet mesh and $0.1608$~km on the meshes of $30{,}000$, $40{,}000$
and $49{,}152$ facets, a difference of $9.4$~m. Any test that holds the
plane fixed must therefore supply it explicitly rather than rely on the
detector returning the same value twice.

The threshold $p \ge 0.02$ is not a fine judgement. Across the bodies examined
the accepted values run from $0.490$ down to $0.091$, and the single rejected
value is $0.0004$: the gap spans more than two orders of magnitude. Where no
candidate qualifies the routine returns no plane, and the neck model is skipped
rather than run on a guess.


